\documentclass[12pt]{article}
\usepackage[T1]{fontenc}
\usepackage[utf8]{inputenc}
\usepackage{lmodern}
\usepackage{textcomp}
\usepackage{enumitem}
\usepackage{geometry}
\geometry{margin=1in}
\begin{document}
\begin{center}
Answer Key - Health Sciences - K-12 Level Assessment 12 \
\small (Answers are ordered exactly as in the question paper.)
\end{center}

\section*{Multiple Choice}
\begin{enumerate}[label=\arabic*.]
\item \textbf{Answer:} A
\\ \textit{Options:} A) Lipid-soluble antioxidant in cell membranes; B) Fast responses to increase calcium absorption with no change in gene expression; C) Regulation of differentiation of adipocytes; D) Regulation of bone turnover
\\ \textit{Note:} The correct answer is "lipid-soluble antioxidant in cell membranes" because vitamin D primarily functions in calcium absorption and bone health regulation, rather than acting as an antioxidant.
\item \textbf{Answer:} D
\\ \textit{Options:} A) is defined as the ratio of VO 2 divided by VCO 2.; B) decreases on a high carbohydrate diet.; C) increases with fasting.; D) goes beyond a value of 1.0 when exogenous carbohydrate is converted to endogenous fat.
\\ \textit{Note:} The respiration quotient (RQ) exceeds 1.0 when exogenous carbohydrates are converted to endogenous fat because this process generates more carbon dioxide than oxygen consumed, indicating a shift in metabolic substrate utilization.
\item \textbf{Answer:} B
\\ \textit{Options:} A) Synthesis of TMP (thymidine monophosphate); B) Decarboxylation of amino acids to form amine neurotransmitters; C) Synthesis of methionine from homocysteine; D) Carboxylation of pyruvate to oxaloacetate
\\ \textit{Note:} Vitamin B$_{6}$ plays a crucial role in the decarboxylation of amino acids, which is essential for the synthesis of amine neurotransmitters that are vital for proper brain function and mood regulation.
\item \textbf{Answer:} B
\\ \textit{Options:} A) Selenophosphate; B) Selenocysteine; C) Selenohistidine; D) Selenate
\\ \textit{Note:} Selenocysteine is the correct answer because it is a unique amino acid that incorporates selenium into proteins, specifically in the synthesis of the 25 human selenoproteins, which play critical roles in various biological processes.
\item \textbf{Answer:} D
\\ \textit{Options:} A) Glucose neurotoxicity; B) Insulin insensitivity; C) Systemic inflammation; D) All of the above
\\ \textit{Note:} The correct answer is "All of the above" because age-related cognitive decline can be influenced by a variety of mechanisms, including neurodegenerative diseases, vascular health issues, and lifestyle factors, all of which contribute to the overall decline in cognitive function as individuals age.
\end{enumerate}

\section*{True / False}
\begin{enumerate}[label=\arabic*.]
\setcounter{enumi}{5}
\item \textbf{Answer:} False
\\ \textit{Options:} A) True; B) False
\\ \textit{Note:} The statement is false because basal metabolic rate is primarily influenced by lean body mass, such as muscle, rather than adipose tissue, which does not contribute significantly to metabolic activity.
\item \textbf{Answer:} True
\\ \textit{Options:} A) True; B) False
\\ \textit{Note:} Biotin is a vital cofactor for the enzyme pyruvate carboxylase, which catalyzes the conversion of pyruvate to oxaloacetate, thereby playing an essential role in gluconeogenesis and energy metabolism.
\item \textbf{Answer:} True
\\ \textit{Options:} A) True; B) False
\\ \textit{Note:} Underwater weighing is considered the most direct method for measuring body composition because it accurately assesses body density by comparing an individual's weight in water to their weight on land, allowing for precise calculations of fat and lean mass.
\item \textbf{Answer:} True
\\ \textit{Options:} A) True; B) False
\\ \textit{Note:} As survival time with a disease increases, more individuals live longer with that condition, leading to a higher number of existing cases in the population, thus increasing prevalence.
\item \textbf{Answer:} False
\\ \textit{Options:} A) True; B) False
\\ \textit{Note:} Prenatal programming can affect all children, regardless of birth weight, as it encompasses the influence of maternal health, nutrition, and environmental factors on fetal development that can lead to long-term health outcomes.
\end{enumerate}

\section*{Long Form Response}
\begin{enumerate}[label=\arabic*.]
\setcounter{enumi}{10}
\item \textbf{Answer:} Proteins play a crucial role in the absorption of cholesterol in the small intestine by facilitating its transport across the intestinal cell membranes. Specific transport proteins, such as Niemann-Pick C$_{1}$-like 1 (NPC 1 L$_{1}$), help cholesterol enter the cells lining the intestine, allowing it to be absorbed into the bloodstream. However, the LDL receptor is not involved in this initial absorption process; instead, it functions later in the body by regulating cholesterol uptake from the blood into cells. The LDL receptor binds to low-density lipoproteins (LDL), which carry cholesterol through the bloodstream, but it does not participate in the direct absorption of dietary cholesterol in the intestine. Thus, while proteins are essential for cholesterol absorption, the LDL receptor plays a different role in cholesterol metabolism after it has been absorbed.
\\ \textit{Note:} correct because it identifies the specific transport protein NPC 1 L$_{1}$ as essential for cholesterol absorption in the small intestine, while clarifying that the LDL receptor's role is limited to regulating cholesterol uptake from the bloodstream after absorption has already occurred.
\item \textbf{Answer:} Osteopenia and osteoporosis are important diagnoses for women over age 80, as they indicate varying levels of bone health. Women diagnosed with osteopenia actually have above average bone mineral density compared to their peers in the same age group. This means that while their bones are weaker than those of younger women, they are still stronger than many of their contemporaries. Understanding these differences is crucial because it helps in managing health risks such as fractures. Early detection and appropriate interventions can support better bone health and overall quality of life for older women.
\\ \textit{Note:} correct because it clarifies that women with osteopenia have better bone mineral density than many of their peers over age 80, which is significant for managing health risks associated with bone health in this age group.
\end{enumerate}

\vfill
\noindent\textit{End of Answer Key}
\end{document}